% ==============================================================================
% Siglas e abreviaturas (pacote acro)
% ------------------------------------------------------------------------------
% Declare aqui cada sigla uma única vez. No texto, use:
%   \ac{chave}   primeira ocorrência: forma extensa seguida da sigla entre
%                parênteses; nas seguintes, só a sigla;
%   \acs{chave}  sempre a sigla (use em títulos e legendas);
%   \acl{chave}  sempre a forma extensa;
%   \acp{chave}  plural da sigla.
% A lista de abreviaturas e siglas é montada automaticamente, em ordem
% alfabética, apenas com as siglas usadas no texto.
%
% Campos:
%   short   sigla;
%   long    forma extensa usada no texto;
%   list    forma extensa na lista (opcional, para ajustar maiúsculas);
%   foreign forma original, quando a sigla vem de outro idioma.
% ==============================================================================

\DeclareAcronym{mkp}{
  short   = MKP,
  long    = problema da mochila multidimensional,
  list    = Problema da mochila multidimensional,
  foreign = Multidimensional Knapsack Problem,
  foreign-babel = english,
}

\DeclareAcronym{ils}{
  short   = ILS,
  long    = busca local iterada,
  list    = Busca local iterada,
  foreign = Iterated Local Search,
  foreign-babel = english,
}

\DeclareAcronym{pli}{
  short = PLI,
  long  = programação linear inteira,
  list  = Programação linear inteira,
}

\DeclareAcronym{ag}{
  short = AG,
  long  = algoritmo genético,
  list  = Algoritmo genético,
}

\DeclareAcronym{ia}{
  short = IA,
  long  = inteligência artificial,
  list  = Inteligência artificial,
}

\DeclareAcronym{abnt}{
  short = ABNT,
  long  = Associação Brasileira de Normas Técnicas,
}
